\subsection{The Helffer--Sjöstrand formula}

In this issue, E denotes a complex Hilbert space. We shall obtain a formula for certain cases of the functional calculus of a partial self-adjoint operator, expressed directly in terms of the resolvent of the operator in question.

We endow R (resp. C) with Lebesgue measure, denoted \(\mu\) and we identify the group R with its dual group by the map \((x,y)\mapsto\exp(2i\pi xy)\) (cf. Corollary 3 of II, p. 236).

For any function \(f\) defined and differentiable on an open set U of \(\mathbf{R}^{2}\) , identified with \(\mathbf{C}\) with real coordinates \(x\) and \(y\) , we denote

\[
{\frac {\partial f}{\partial \overline {{{z}}}}} = {\frac {1}{2}} {\Bigl (} {\frac {\partial f}{\partial x}} + i {\frac {\partial f}{\partial y}} {\Bigr)}
\]

(cf. VAR, R2, 8.8.10, p. 24).

Lemma 5. --- Let \(f \in \mathcal{D}(\mathbf{R})\) . There exists a function \(\tilde{f}\) in \(\mathcal{D}(\mathbf{C})\) which coincides with \(f\) on \(\mathbf{R}\) and satisfies

\[
\frac {\partial \widetilde {f}}{\partial \overline {{z}}} (x, 0) = 0\tag{11}
\]

for all \(x \in \mathbf{R}\) . We then have

\[
\frac {\partial \widetilde {f}}{\partial y} (x, 0) = i f ^ {\prime} (x)\tag{12}
\]

for all \(x \in R\) and there exists a real number C ≥ 0 such that

\[
\left| \frac {\partial \widetilde {f}}{\partial \overline {{z}}} (x, y) \right| \leqslant C | y |\tag{13}
\]

for all \((x,y)\in \mathbf{R}^2\)

There exists \(\varphi\in\mathcal{D}(\mathbf{R})\) whose support is contained in \([-2,2]\) and which is equal to 1 on \([-1,1]\) (lemma 1 of IV, p. 196). Set

\[
\widetilde {f} (x, y) = \bigl (f (x) + i y f ^ {\prime} (x) \bigr) \varphi (y)
\]

for \((x,y)\in\mathbf{R}^{2}\) . We have \(\widetilde{f}\in\mathcal{D}(\mathbf{C})\) and \(\widetilde{f}\) coincides with f on R. Moreover, for any \((x,y)\in\mathbf{R}^{2}\) , it follows that

\[
\frac {\partial \widetilde {f}}{\partial \overline {{z}}} (x, y) = \frac {1}{2} \Big ((i f (x) - y f ^ {\prime} (x)) \varphi^ {\prime} (y) + i y f ^ {\prime \prime} (x) \varphi (y) \Big).
\]

As the function \(\varphi\) is equal to 1 in a neighborhood of 0, we have \(\varphi'(0) = 0\) whence (11).

Let \(\tilde{f}\) be in \(\mathcal{D}(\mathbf{C})\) satisfying (11). Formula (12) follows from this, and the bound (13) is obtained using the mean value theorem (FVR, I, p. 23, th. 2).

We say that \(\tilde{f}\) is an almost analytic extension of f.

Lemma 6. --- Let \(\varepsilon > 0\) . The function \(\sigma_{\varepsilon}\) defined on \(\mathbf{R}\) by

\[
\sigma_ {\varepsilon} (x) = \frac {2 i \varepsilon x}{x ^ {2} + \varepsilon^ {2}}
\]

belongs to \(L^{2}(\mathbf{R})\) Its Fourier transform is the class in \(L^{2}(\mathbf{R})\) of the function \(\eta_{\varepsilon}\) which vanishes at 0 and satisfies

\[
\eta_ {\varepsilon} (y) = \frac {2 \pi \varepsilon y}{| y |} e ^ {- 2 \pi \varepsilon | y |}
\]

for all \(y \neq 0\) .

One has \(\sigma_{\varepsilon} \in \mathrm{L}^{2}(\mathbf{R})\) by prop. 3 of IV, p. 199, and the class of the function \(\eta_{\varepsilon}\) belongs to \(\mathrm{L}^{2}(\mathbf{R}) \cap \mathrm{L}^{1}(\mathbf{R})\) . For every \(x \in \mathbf{R}\) , we have

\[
\begin{array}{r} \overline {{\mathcal {F}}} (\eta_ {\varepsilon}) (x) = 2 \pi \varepsilon \int_ {\mathbf {R} _ {+}} e ^ {2 \pi (i x - \varepsilon) y} d y - 2 \pi \varepsilon \int_ {\mathbf {R} _ {-}} e ^ {2 \pi (i x + \varepsilon) y} d y \\ = \frac {\varepsilon}{\varepsilon - i x} - \frac {\varepsilon}{\varepsilon + i x} = \frac {2 i \varepsilon x}{x ^ {2} + \varepsilon^ {2}}, \end{array}
\]

whence \(\overline{\mathcal{F}}(\eta_{\varepsilon}) = \sigma_{\varepsilon}\) The result then follows from the Fourier inversion formula in \(\mathrm{L}^2(\mathbf{R})\) (corollary of Theorem 2 of II, p. 220).

Lemma 7. --- Let \(f \in \mathcal{D}(\mathbf{R})\) and let \(\widetilde{f} \in \mathcal{D}(\mathbf{C})\) be an almost analytic extension of \(f\) . For \(\varepsilon > 0\) , define \(f_{\varepsilon}\) on \(\mathbf{R}\) by

\[
f _ {\varepsilon} (x) = - \frac {1}{2 i \pi} \int_ {\mathbf {R}} \biggl (\frac {\widetilde {f} (y + i \varepsilon)}{y - x + i \varepsilon} - \frac {\widetilde {f} (y - i \varepsilon)}{y - x - i \varepsilon} \biggr) d y.
\]

Then \(f_{\varepsilon}\) is continuous and bounded, and \(f_{\varepsilon}\) converges to \(f\) in \(\mathcal{C}_b(\mathbf{R})\) as \(\varepsilon\) tends to 0.

The continuity of \(f_{\varepsilon}\) is a consequence of INT, IV, p. 144, § 4, n° 3, cor. 1, since \(\widetilde{f}\) is compactly supported. Moreover, if \(r\) is such that the support of \(\widetilde{f}\) is contained in \([-r, r] \times \mathbf{R}\) , we have

\[
| f _ {\varepsilon} (x) | \leqslant 2 \| \widetilde {f} \| _ {\infty} \int_ {- r} ^ {r} \frac {1}{\sqrt {(x - y) ^ {2} + \varepsilon^ {2}}} d y \leqslant \frac {4 r}{\varepsilon} \| \widetilde {f} \| _ {\infty},
\]

hence \(f_{\varepsilon}\) is bounded.

The first-order Taylor expansion (FVR, I, p. 30, prop. 3) and formula (13) show that there exist M ≥ 0 and a function \(\varrho_{1}\) with values in \(R^{2}\) such that

\[
\widetilde {f} (y + i \gamma) = f (y) + i \gamma f ^ {\prime} (y) + \gamma^ {2} \varrho_ {1} (y; \gamma), \text {   and   } | \varrho_ {1} (y; \gamma) | \leqslant M,
\]

for all \(y \in \mathbf{R}\) and \(\gamma \in \mathbf{R}\) . Since the map \(y \mapsto \widetilde{f}(y + i\gamma)\) has support contained in \([-r, r]\) for all \(\gamma \in \mathbf{R}\) , the map \(y \mapsto \varrho_{1}(y; \gamma)\) has support contained in \([-r, r]\) , for all \(\gamma \in \mathbf{R}\) .

Let \(g_{\varepsilon}\) the function defined on \(\mathbf{R}^2\) by

\[
g _ {\varepsilon} (x, y) = \frac {\widetilde {f} (y + i \varepsilon)}{y - x + i \varepsilon} - \frac {\widetilde {f} (y - i \varepsilon)}{y - x - i \varepsilon}.
\]

For every \((x,y)\in \mathbf{R}^2\) and \(\varepsilon >0\) , we obtain

\[
\begin{array}{r l} g _ {\varepsilon} (x, y) = - \frac {2 i \varepsilon}{(x - y) ^ {2} + \varepsilon^ {2}} f (y) + \frac {2 i \varepsilon (y - x)}{(x - y) ^ {2} + \varepsilon^ {2}} f ^ {\prime} (y) \\ & \quad + \varepsilon^ {2} \bigg (\frac {\varrho_ {1} (y ; \varepsilon)}{y - x + i \varepsilon} - \frac {\varrho_ {1} (y ; - \varepsilon)}{y - x - i \varepsilon} \bigg). \end{array}
\]

Let \(x \in \mathbf{R}\) and \(\varepsilon > 0\) . We have

\[
\begin{array}{c} \left| \varepsilon^ {2} \int_ {\mathbf {R}} \left(\frac {\varrho_ {1} (y ; \varepsilon)}{y - x + i \varepsilon} - \frac {\varrho_ {1} (y ; - \varepsilon)}{y - x - i \varepsilon}\right) d y \right| \leqslant 2 \mathrm{M} \varepsilon^ {2} \int_ {- r} ^ {r} \frac {d y}{\sqrt {(x - y) ^ {2} + \varepsilon^ {2}}} \\ \leqslant 4 \mathrm{Mr} \varepsilon . \end{array}
\]

Consequently, for every \(x \in \mathbf{R}\) , it follows that

\[
f _ {\varepsilon} (x) = - \frac {1}{2 i \pi} \int_ {\mathbf {R}} g _ {\varepsilon} (x, y) d y = (f * \delta_ {\varepsilon}) (x) - \frac {1}{2 i \pi} (f ^ {\prime} * \sigma_ {\varepsilon}) (x) + k _ {\varepsilon} (x)
\]

where \(\delta_{\varepsilon}\) and \(\sigma_{\varepsilon}\) are the functions on \(\mathbf{R}\) defined by

\[
\delta_ {\varepsilon} (x) = \frac {1}{\pi} \frac {\varepsilon}{x ^ {2} + \varepsilon^ {2}}, \qquad \sigma_ {\varepsilon} (x) = \frac {2 i \varepsilon x}{x ^ {2} + \varepsilon^ {2}},
\]

and \(\| k_{\varepsilon}\|_{\infty}\leqslant 2\mathrm{Mr}\varepsilon .\)

The function \(\mathcal{F}(f') \in \mathcal{S}(\mathbf{R})\) is integrable (prop. 18 of IV, p. 219 and prop. 13 of IV, p. 213). By lemma 6, it follows that

\[
\begin{array}{r l r} & & {\int_ {\mathbf {R}} | \mathcal {F} (f ^ {\prime}) (y) \mathcal {F} (\sigma_ {\varepsilon}) (y) | d y = 2 \pi \varepsilon \int_ {\mathbf {R}} | \mathcal {F} (f ^ {\prime}) (y) | e ^ {- 2 \pi \varepsilon | y |} d y} \\ & & {\leqslant 2 \pi \varepsilon \int_ {\mathbf {R}} | \mathcal {F} (f ^ {\prime}) (y) | d y,} \end{array}
\]

hence \(\mathcal{F}(f')\mathcal{F}(\sigma_{\varepsilon})\) converges to 0 in \(\mathrm{L}^{1}(\mathbf{R})\) as \(\varepsilon\) tends to 0. Since \(f'\) and \(\sigma_{\varepsilon}\) belong to \(\mathrm{L}^{2}(\mathbf{R})\) , prop. 14 of II, p. 223 implies that \(f'*\sigma_{\varepsilon}=\overline{\mathcal{F}}(\mathcal{F}(f')\mathcal{F}(\sigma_{\varepsilon}))\) converges to 0 in \(\mathcal{C}_{b}(\mathbf{R})\) .

For every \(\varepsilon > 0\) , the positive measure \(\delta_{\varepsilon} \cdot dx\) on \(\mathbf{R}\) has total mass 1 (cf. FVR, III, p. 7). The set of measures \(\delta_{\varepsilon} \cdot dx\) for \(\varepsilon > 0\) and the filter induced on this set by the filter of neighborhoods of 0 in \(\mathbf{R}_{+}^{*}\) satisfy the hypotheses of lemma 4 of INT, VIII, p. 137, § 2, no. 7. The measures \(\delta_{\varepsilon} \cdot dx\) therefore converge in \(\mathcal{M}^{1}(\mathbf{R})\) toward the point measure \(\varepsilon_{0}\) as \(\varepsilon\) tends to 0. It follows that \(f * \delta_{\varepsilon} \to f\) in \(\mathcal{C}_{b}(\mathbf{R})\) (INT, VIII, p. 163, § 4, no. 4). The lemma is proved.

We denote by \(\mu\) the Lebesgue measure on \(\mathbf{C}\).

THEOREM 2 (Helffer--Sjöstrand). --- Let u be a partial self-adjoint operator on E. Let \(f \in \mathcal{D}(\mathbf{R})\) and \(\widetilde{f} \in \mathcal{D}(\mathbf{C})\) be an almost analytic extension of \(f\) . Let h be the map from \(\mathbf{C}\) into \(\mathcal{L}(\mathrm{E})\) defined by

\[
h (\lambda) = \frac {\partial \widetilde {f}}{\partial \overline {{z}}} (\lambda) \mathrm{R} (u, \lambda)
\]

if \(\lambda\in\mathbf{C} - \mathbf{R}\) and \(h(\lambda)=0\) if \(\lambda\in\mathbf{R}\) . Then h is \(\mu\) -integrable on \(\mathbf{C}\) and

\[
f (u) = - \frac {1}{\pi} \int_ {\mathbf {C}} h (\lambda) d \mu (\lambda).
\]

The map h is measurable and its support is compact. Since u is self-adjoint, we have \(\|\mathrm{R}(u,\lambda)\| \leqslant |\mathcal{I}(\lambda)|^{-1}\) for all \(\lambda \in \mathbf{C} - \mathbf{R}\) (prop. 17 of IV, p. 248). The map h is therefore bounded according to formula (13), and consequently it is integrable on \(\mathbf{C}\).

Let \(\varepsilon \in \mathbf{R}_{+}^{*}\) . We denote \(\mathrm{F}_{\varepsilon}^{+}\) (resp. \(\mathrm{F}_{\varepsilon}^{-}\) ) the closed set in \(\mathbf{C}\) of the \(\lambda \in \mathbf{C}\) such that \(\mathcal{I}(\lambda) \geqslant \varepsilon\) (resp. \(\mathcal{I}(\lambda) \leqslant -\varepsilon\) ). We have

\[
\int_ {\mathbf {C}} h (\lambda) d \mu (\lambda) = \lim _ {\varepsilon \to 0} \Bigl (\int_ {\mathrm{F} _ {\varepsilon} ^ {+}} h (\lambda) d \mu (\lambda) + \int_ {\mathrm{F} _ {\varepsilon} ^ {-}} h (\lambda) d \mu (\lambda) \Bigr).
\]

Let r > 0 such that the support of h is contained in \(\mathbf{C} = [-r, r]^{2}\) . Denote by \(R_{\varepsilon}^{+}\) the box \([-2r, 2r] \times [\varepsilon, 2r]\) in \(\mathbf{C}\). It is a locally polyhedral subset of \(\mathbf{C}\) (VAR, R2, 11.3, p. 48). It satisfies the following conditions:

\begin{enumerate}
\def\labelenumi{(\roman{enumi})}
\item
  One has \(R_{\varepsilon}^{+} \subset 2C \cap F_{\varepsilon}^{+}\) ;
\item
  The set \(R_{\varepsilon}^{+}\) contains the intersection of \(F_{\varepsilon}^{+}\) and the support of h;
\item
  The regular boundary \(\partial R_{\varepsilon}^{+}\) (VAR, R2, 11.3.2, p. 49) contains the segment \(S_{\varepsilon} = [-r, r] + i\varepsilon \subset C\) ;
\item
  One has \(h(\lambda) = 0\) if \(\lambda \in \partial \mathrm{R}_{\varepsilon}^{+} - \mathrm{S}_{\varepsilon}\) .
\end{enumerate}

Denote by \(d\lambda\) (resp. \(d\overline{\lambda}\) ) the differential 1-form on \(\mathbf{C}\) of the identity map of \(\mathbf{C}\) (resp. of complex conjugation). Denote by \(g\) the function on \(\mathbf{C} - \mathbf{R}\) with values in \(\mathcal{L}(\mathrm{E})\) such that \(g(\lambda) = \widetilde{f}(\lambda)\mathrm{R}(u,\lambda)\) for \(\lambda \in \mathbf{C} - \mathbf{R}\) . Let \(\omega = g d\lambda\) ; it is a differential form of degree 1 on \(\mathbf{C} - \mathbf{R}\) with compact support and values in \(\mathcal{L}(\mathrm{E})\) . Since the resolvent of \(u\) is holomorphic (prop. 14 of IV, p. 246), we have

\[
d \omega = \left(\frac {\partial \widetilde {f}}{\partial \overline {{z}}} (\lambda) \mathrm{R} (u, \lambda) + \widetilde {f} (\lambda) \frac {\partial}{\partial \overline {{z}}} \mathrm{R} (u, \lambda)\right) d \overline {{\lambda}} \wedge d \lambda = - h (\lambda) d \lambda \wedge d \overline {{\lambda}}.
\]

The vector measure associated with \(d\omega\) (VAR, R2, 10.4.3, p. 43) is the density measure -2ih with respect to Lebesgue measure. Applying Stokes' formula to the locally polyhedral set \(R_{\varepsilon}^{+}\) and to the differential form \(\omega\) (VAR, R2, 11.3.4, p. 49), we thus obtain

\[
\frac {i}{2} \int_ {\mathrm{R} _ {\varepsilon} ^ {+}} d \omega = \frac {i}{2} \int_ {\partial \mathrm{R} _ {\varepsilon} ^ {+}} \omega = \frac {i}{2} \int_ {\mathrm{S} _ {\varepsilon}} \omega = \frac {i}{2} \int_ {- r} ^ {r} \widetilde {f} (y + i \varepsilon) \mathrm{R} (u, y + i \varepsilon) d y,
\]

whence

\[
\int_ {\mathrm{F} _ {\varepsilon} ^ {+}} h (\lambda) d \mu (\lambda) = - \frac {i}{2} \int_ {\mathbf {R}} \widetilde {f} (y + i \varepsilon) \mathrm{R} (u, y + i \varepsilon) d y.
\]

Reasoning in the same way for \(F_{\varepsilon}^{-}\) , we obtain

\[
\int_ {\mathrm{F} _ {\varepsilon} ^ {-}} h (\lambda) d \mu (\lambda) = - \frac {i}{2} \int_ {\mathbf {R}} \widetilde {f} (y - i \varepsilon) \mathrm{R} (u, y - i \varepsilon) d y,
\]

and we conclude that the integral of \(h\) over \(\mathbf{C}\) is the limit when \(\varepsilon\to0\) of

\[
v _ {\varepsilon} = \frac {i}{2} \int_ {\mathbf {R}} \Bigl (\widetilde {f} (y + i \varepsilon) \mathrm{R} (u, y + i \varepsilon) - \widetilde {f} (y - i \varepsilon) \mathrm{R} (u, y - i \varepsilon) \Bigr) d y.
\]

By prop. 7, we have \(v_{\varepsilon} = \pi f_{\varepsilon}(u)\) , where \(f_{\varepsilon}\) is the function defined on \(\mathbf{R}\) by

\[
f _ {\varepsilon} (x) = - \frac {1}{2 i \pi} \int_ {\mathbf {R}} \biggl (\frac {\widetilde {f} (y + i \varepsilon)}{y - x + i \varepsilon} - \frac {\widetilde {f} (y - i \varepsilon)}{y - x - i \varepsilon} \biggr) d y.
\]

Since \(f_{\varepsilon} \to f\) as \(\varepsilon \to 0\) in \(\mathcal{C}_{b}(\mathbf{R})\) (Lemma 7), the endomorphism \(v_{\varepsilon} = \pi f_{\varepsilon}(u)\) converges to \(\pi f(u)\) in \(\mathcal{L}(\mathrm{E})\) (prop. 5 of IV, p. 275). The theorem is proved.

